%=====================================================================
% 引言（Introduction）
% 原文：P. Etingof, arXiv:2409.03117v1, p.4–6
%=====================================================================
\section*{引言}
\phantomsection
\addcontentsline{toc}{section}{引言}

\begin{keytakeaways}
\begin{itemize}
  \item \textbf{物理图像：} 量子场论与弦论在近四十年的发展，把代数拓扑、范畴论、微分与代数几何、表示论、
        组合学乃至数论引入了物理学；本书试图以数学家能接受的方式，梳理量子场论的基本框架。
  \item \textbf{数学工具：} 泛函积分与概率测度（Wiener 积分）、最速下降与驻相分析、Feynman 演算与组合计数、
        矩阵积分与胖图、算子的自伴扩张与半群、超几何与 Berezin 积分、Wightman 公理、共形对称与 Virasoro 代数。
  \item \textbf{关键结论：} 全书以“维数递增”为线索：$0$ 维（矩阵积分与组合）$\to$ $0+1$ 维（量子力学）$\to$
        $d+1$ 维（自由场与微扰论）$\to$ $1+1$ 维（共形场论）；每一层都展示同一套形式计算（路径积分、Feynman 图）
        如何被翻译成严格的数学对象。
\end{itemize}
\end{keytakeaways}

\bigskip
\noindent\rule{\textwidth}{0.6pt}
\bigskip

物理学一直是数学若干核心领域（分析、微分方程、概率论等）的重要动机与应用来源。然而，过去四十年间
量子场论与弦论的发展，把这两个学科之间的互动推到了前所未有的高度：代数拓扑、范畴论、微分几何与代数几何、
表示论、组合学、乃至数论这些传统上“纯粹”的数学领域，都被纳入了物理学。
这种互动在两个方向上都极富成果，也使得物理学家有必要了解现代数学的基础，数学家有必要了解现代物理学的基础。
物理学家学得很快，如今优秀的物理学家对相关数学领域的理解，往往与职业数学家一样深刻。另一方面，许多数学家
却一直裹足不前，被物理文献缺乏严格性所阻，更被其不同的表述方式所阻。特别是，量子场论的基本框架——
哪怕只是为了理解其更高级（也更有数学趣味）的部分所必需的框架——对数学家来说已经相当陌生。
尽管如此，量子场论的许多基本思想实际上可以用严格且数学上自然的方式表述出来。这正是本书的主要目标。

本书是为量子场论基本数学结构的研究生课程所写的、略作扩充的讲义；该课程由作者于 2002 年在 MIT 数学系讲授，
2023 年又讲授了一次。读者不应指望从本书学会量子场论——这是不可能的：例如，我对这门学科的了解就少于一个
初学物理的研究生。更确切地说，正如上文所述，本书旨在以有数学动机的方式呈现量子场论的基本框架，突出它与
数学各领域的联系。因此，它可以为读者阅读这一类型的更高级著作（例如 \cite{QFS}）作准备，或者为读者从数学家的
视角阅读通常的量子场论教科书或讲义（例如 \cite{Co,W,IZ,PS}）作准备。请注意，许多重要的材料都包含在习题中，
我强烈建议读者在阅读时解答这些习题。

我们从经典力学、量子力学与场论的一般讨论开始（第 1 章）。接着证明经典渐近分析中的最速下降公式与驻相公式，
它们为路径积分的微扰计算提供了一个有限维模型（第 2 章）。然后，在第 3 章中我们发展 Feynman 演算——
微扰量子场论中主要的组合工具。为说明 Feynman 演算，我们给出它在计数组合学中的若干应用
（矩阵树定理及其特例）。

在第 4 章中，我们把 Feynman 演算推广到矩阵积分，并证明：对这类积分，Feynman 图被胖图（曲面）所取代，
从而 $1/N$ 渐近展开（$N$ 为矩阵大小）的系数是给定亏格的胖图之和。这使得我们能够在第 5 章证明
Harer–Zagier 关于曲线模空间欧拉示性数的定理，并在第 6 章得到平面图的非平凡计数。

然而，以上全部材料讨论的都是 $0$ 维时空（或者开玩笑地说，$-1$ 维空间）中的量子场论。要联系真实物理，
我们必须至少升高一维，即考虑 $0+1$ 维时空（或 $0$ 维空间）中的量子场论，这就是量子力学（第 7 章）。
我们先回顾经典力学的拉格朗日形式（拉格朗日量、最小作用量原理），然后对这一形式作量子化，发展量子力学的
路径积分方法。也就是说，我们用 Feynman 图描述量子力学路径积分的微扰展开，并给出若干例子。我们还说明：
通过时间的 Wick 转动 $t\mapsto it$，量子力学路径积分与随机过程理论中（严格定义的）Wiener 积分相联系。

在第 8 章中，我们在回顾经典力学的哈密顿形式之后，描述它的量子化；这为非微扰量子力学提供了严格基础。
我们还证明 Feynman–Kac 公式，它把拉格朗日途径与哈密顿途径得到的关联函数联系起来。

在第 9、10 章中，我们讨论前几章材料的超推广，即描述费米子的经典力学与量子力学。我们先回顾超几何与
超流形上的 Berezin 积分理论，进而把 Feynman 演算推广到超情形。然后讨论费米子的经典力学与量子力学。

在第 11 章中，我们终于进入真正的量子场论：$d+1$ 维时空，$d\ge 1$。我们先回顾拉格朗日经典场论，
然后过渡到它的量子化，特别是描述自由玻色子与自由费米子的理论。我们回顾旋量的经典理论
（特别是实旋量的 Bott 周期性，模 $8$），并用它来描述费米子拉格朗日量中可能的动能项与质量项。
接着转向经典场论与量子场论中的哈密顿形式，并讨论 Wightman 公理。最后从这一观点描述自由标量玻色子的量子理论。

在第 12 章中，我们描述微扰重整化理论的基础。特别地，我们讨论 Feynman 振幅的紫外发散，以及通过在
拉格朗日量中引入依赖动量空间截断 $\Lambda$ 的抵消项来正规化这类发散。我们定义超可重整、可重整与不可重整
理论，各种理论及拉格朗日量中各项的临界维数，并讨论关键的例子。

最后，在第 13 章中我们简要介绍二维共形场论。在回顾 $1+1$ 维无质量标量场的经典场论之后，我们将其量子化，
并由无穷维 Heisenberg 李代数的 Fock 表示构造其 Hilbert 空间。我们证明该理论的配分函数（用 $\zeta$ 函数
正规化归一化）是模不变的，这反映了它的共形对称性。然后我们证明该理论的 Hilbert 空间承载着
$\Bbb C^\times$ 上多项式向量场李代数 $W$ 的两个相互对易的射影作用，它表达了无穷小共形对称性。
我们说明这一作用确实是射影的，即 $W$ 的两个副本都被其非平凡的中心扩张——Virasoro 代数（共形反常）——
所取代。接着我们讨论该理论的圆值版本、顶点算子与 $T$ 对偶。我们还简要讨论 $1+1$ 维自由费米子的量子理论
以及 Wess–Zumino–Witten 模型。

{\bf 致谢。}我非常感谢 \cite{QFS} 的合作者；若不是我们近三十年前在这个项目上的合作，本书肯定不会问世。
特别感谢 David Kazhdan，是他促使我学习量子场论的基础；感谢 Edward Witten，我关于这门学科所知道的东西
几乎全部来自他；感谢 Dan Freed，没有他细致的笔记与解释，本书将不可能完成；感谢 Pierre Deligne，
他以深刻的数学洞见、清晰与优雅，极大地促进并深化了我们的学习。我也感谢 2002 年与 2023 年 MIT 课程的
参与者，正是这些课程催生了本书。本工作部分受 NSF 基金 DMS-2001318 支持。

% [Tutor Note: 原文此处作 "a number of ins applications"，其中 "ins" 当为 "nice" 之笔误（原文如此）；
%  译文按“若干应用”处理，特此说明，以免读者误以为漏译。]
% [Tutor Note: 作者自称“我对这门学科的了解少于一个初学物理的研究生”，是讲义的自我定位而非字面事实；
%  本书的价值不在物理直觉的训练，而在把路径积分/Feynman 图当作形式幂级数与组合对象来处理，
%  读者应始终注意“形式计算”与“真实物理时空”之间的落差。]
% [Tutor Note: 引言把全书线索概括为“维数递增”：$0$ 维矩阵积分 → $0+1$ 维量子力学 → $d+1$ 维自由场
%  与微扰论 → $1+1$ 维共形场论。阅读时建议对照目录抓住这条主线，否则容易把第 3–6 章误当作纯组合学。]
% [Tutor Note: 第 13 章所谓“两个相互对易的射影作用”应理解为 Witt 代数 $\mathrm{Vect}(S^1)$ 的中心扩张——
%  Virasoro 代数；中心荷的数值既依赖于所取的正规序约定，也依赖于是否计入零模。建议对照
%  Kac, "Infinite Dimensional Lie Algebras" 第 2 章，以及 Di Francesco–Mathieu–Sénéchal 第 5 章。]
